\subsection{Разработка MBeans}

В соответствии с вариантом задания были разработаны два MBean для мониторинга статистики проверки точек.

\subsubsection{PointCounterMBean}

Первый MBean отвечает за подсчёт общего количества установленных пользователем точек и количества точек, не попавших в область (промахов). При достижении числа точек, кратного 15, MBean отправляет соответствующее уведомление. Учитываются как клики по графику, так и отправка формы.

\begin{lstlisting}[caption=Листинг интерфейса \texttt{PointCounterMBean}.]
package weblab.jmx;

public interface PointCounterMBean {
    int getTotalPoints();

    int getMissPoints();

    void resetStats();

    void addPointEvent(boolean hit);
}
\end{lstlisting}

Класс-реализация поддерживает отправку JMX-уведомлений через реализацию интерфейса \texttt{NotificationBroadcaster}.

\begin{lstlisting}[caption=Листинг реализации \texttt{PointCounter}.]
package weblab.jmx;

import javax.management.*;
import java.util.concurrent.atomic.AtomicInteger;

public class PointCounter implements PointCounterMBean, NotificationBroadcaster {
    private final AtomicInteger totalPoints = new AtomicInteger(0);
    private final AtomicInteger missPoints = new AtomicInteger(0);
    private long notificationSequence = 1;
    private final NotificationBroadcasterSupport notificationSupport = new NotificationBroadcasterSupport();

    @Override
    public int getTotalPoints() {
        return totalPoints.get();
    }

    @Override
    public int getMissPoints() {
        return missPoints.get();
    }

    @Override
    public void resetStats() {
        totalPoints.set(0);
        missPoints.set(0);
    }

    @Override
    public void addPointEvent(boolean hit) {
        int newTotal = totalPoints.incrementAndGet();
        if (!hit) missPoints.incrementAndGet();

        if (newTotal % 15 == 0) sendNotification(newTotal, hit);
    }

    private void sendNotification(int total, boolean lastHit) {
        Notification notification = new Notification(
                "point.count.multiple.15",
                this,
                notificationSequence++,
                System.currentTimeMillis(),
                String.format("Total points count = %d (last hit = %s)", total, lastHit)
        );
        notificationSupport.sendNotification(notification);
    }

    @Override
    public void addNotificationListener(NotificationListener listener, NotificationFilter filter, Object handback) throws IllegalArgumentException {
        notificationSupport.addNotificationListener(listener, filter, handback);
    }

    @Override
    public void removeNotificationListener(NotificationListener listener) throws ListenerNotFoundException {
        notificationSupport.removeNotificationListener(listener);
    }

    @Override
    public MBeanNotificationInfo[] getNotificationInfo() {
        return new MBeanNotificationInfo[]{
                new MBeanNotificationInfo(
                        new String[]{"point.count.multiple.15"},
                        Notification.class.getName(),
                        "Sent when total point count becomes a multiple of 15"
                )
        };
    }
}
\end{lstlisting}

\subsubsection{MissPercentageMBean}

Второй MBean вычисляет процентное отношение числа промахов к общему количеству кликов пользователя по координатной плоскости.

\begin{lstlisting}[caption=Листинг интерфейса \texttt{MissPercentageMBean}.]
package weblab.jmx;

public interface MissPercentageMBean {
    double getMissPercentage();

    String getMissPercentageFormatted();

    void resetStats();

    void addClickEvent(boolean hit);
}
\end{lstlisting}

\begin{lstlisting}[caption=Листинг реализации \texttt{MissPercentage}.]
package weblab.jmx;

import java.text.DecimalFormat;
import java.util.concurrent.atomic.AtomicInteger;

public class MissPercentage implements MissPercentageMBean {
    private final AtomicInteger totalClicks = new AtomicInteger(0);
    private final AtomicInteger missClicks = new AtomicInteger(0);
    private final DecimalFormat df = new DecimalFormat("0.000");

    @Override
    public double getMissPercentage() {
        int total = totalClicks.get();
        if (total == 0) return 0.0;

        return (missClicks.get() * 100.0) / total;
    }

    @Override
    public String getMissPercentageFormatted() {
        return df.format(getMissPercentage()) + "%";
    }

    @Override
    public void resetStats() {
        totalClicks.set(0);
        missClicks.set(0);
    }

    @Override
    public void addClickEvent(boolean hit) {
        totalClicks.incrementAndGet();
        if (!hit) missClicks.incrementAndGet();
    }
}
\end{lstlisting}

\subsubsection{Регистрация MBean}

Регистрация обоих MBean выполняется в классе \texttt{JMXConfig} с областью видимости \texttt{@ApplicationScoped}, что гарантирует существование единственного экземпляра на всё приложение.

\begin{lstlisting}[caption=Регистрация \texttt{MBean}.]
package weblab.jmx;

import jakarta.annotation.PostConstruct;
import jakarta.enterprise.context.ApplicationScoped;

import javax.management.MBeanServer;
import javax.management.ObjectName;
import java.lang.management.ManagementFactory;

@ApplicationScoped
public class JMXConfig {
    private PointCounterMBean pointCounter;
    private MissPercentageMBean missPercentage;

    @PostConstruct
    public void init() {
        try {
            MBeanServer mbs = ManagementFactory.getPlatformMBeanServer();

            pointCounter = new PointCounter();
            mbs.registerMBean(pointCounter, new ObjectName("weblab:type=PointCounter"));

            missPercentage = new MissPercentage();
            mbs.registerMBean(missPercentage, new ObjectName("weblab:type=MissPercentage"));

            System.out.println("JMX MBeans registered successfully");
        } catch (Exception e) {
            System.err.println("Failed to register JMX MBeans");
        }
    }

    public PointCounterMBean getPointCounter() {
        return pointCounter;
    }

    public MissPercentageMBean getMissPercentage() {
        return missPercentage;
    }
}
\end{lstlisting}

\subsubsection{Интеграция с приложением}

Вызов методов MBean производится в бине \texttt{PointCheckBean} при обработке кликов по графику (метод \texttt{submitFromGraph()}), а также отправки формы (метод \texttt{submitForm(int y)}). При этом первый MBean получает информацию обо всех установленных точках, а второй — только о кликах по координатной плоскости.

\begin{lstlisting}[caption=Фрагмент интеграции \texttt{MBean}.]
@Inject
private JMXConfig jmxConfig;

public void submitFromGraph() {
    // ... smth cropped ...

    boolean hit = results.get(0).hit();
    
    jmxConfig.getPointCounter().addPointEvent(hit);
    jmxConfig.getMissPercentage().addClickEvent(hit);

    // ... smth cropped ...
}

public void submitForm(int y) {
    // ... smth cropped ...

    for (CheckResult result : results) {
        jmxConfig.getPointCounter().addPointEvent(result.hit());
    }

    // ... smth cropped ...
}
\end{lstlisting}